\label{sec:responsable}

Los resultados de este trabajo son, en sí mismos, el argumento de por qué un
detector de fauna no debe desplegarse como oráculo. Tres de ellos lo dicen con
números propios.

\emph{El desempeño se derrumba fuera de las cámaras conocidas.} La brecha entre
ubicaciones vistas y nuevas (Sección~\ref{sec:generalizacion}) significa que
validar un modelo en las cámaras donde se entrenó sobreestima lo que hará en una
cámara nueva. Un proyecto que instala un sensor en un sitio no muestreado y
confía en la cifra de validación está usando un número que no corresponde a su
situación. La consecuencia operativa es concreta: las métricas que se publican
junto a un modelo de este dominio deben declarar en qué ubicaciones se midieron.

\emph{El error no es uniforme entre especies, y es peor donde más importa.} Las
especies con menos datos de entrenamiento son las que el modelo detecta peor, y
suelen ser también las de mayor interés para conservación, las raras. Los dos
errores no cuestan lo mismo en este dominio: un falso negativo, no registrar una
especie rara que sí estaba, puede hacer concluir que un sitio está vacío y
justificar no protegerlo, mientras que un falso positivo cuesta sobre todo tiempo
de verificación humana; el daño se concentra justo donde el modelo es más débil.
Un reporte automático subestima sistemáticamente a las especies cuya presencia o
ausencia puede motivar una decisión de manejo. Reportar un único mAP agregado
oculta ese sesgo; por eso el desglose por especie con su soporte al lado no es un
detalle de presentación sino una condición para que el número sea utilizable.

\emph{El modelo no mira donde uno supone.} La razón de masa de Grad-CAM
(Sección~\ref{sec:gradcam}) muestra que buena parte de la evidencia que el modelo
usa no cae sobre el animal. Un sistema que acierta apoyándose en el fondo acierta
hasta que el fondo cambia, y no hay nada en su salida que avise cuándo eso
ocurrió. La confianza que el detector reporta no es una medida de que esté
mirando lo correcto.

De ahí, cuatro compromisos para un despliegue responsable, respaldados por las
mediciones de este trabajo.

\emph{Asistir, no decidir.} El uso defendible es el que MegaDetector adoptó
explícitamente~\cite{beery2019efficient}: filtrar el volumen de imágenes vacías,
que es el trabajo tedioso, y dejar la identificación de especie a revisión
humana, especialmente en las especies raras, donde el modelo es menos fiable y el
costo de un error es mayor.

\emph{Declarar la incertidumbre junto al número.} Un conteo de eventos sin su
error invita a conclusiones que no aguanta; por eso el módulo de monitoreo
reporta su error especie por especie contra las anotaciones humanas, en vez de
dejar que el lector lo suponga.

\emph{Cuidar el dato de ubicación.} Las coordenadas de las cámaras trampa son
información sensible: un registro de dónde y cuándo aparece una especie amenazada
es lo que necesita quien quiera cazarla. Un pipeline que indexa y publica
detecciones amplifica ese riesgo a escala, así que lo que se comparte debe ir
agregado o con la ubicación degradada deliberadamente.

\emph{No confundir ausencia de detección con ausencia de animal.} La ocupación
que el módulo del entregable \#5 reporta es \emph{naive}: una especie presente
pero no fotografiada cuenta como ausente. Corregirlo exige modelar la
probabilidad de detección~\cite{mackenzie2002estimating}, y mientras no se haga,
la cifra es una cota inferior, no una estimación. Decirlo es parte del reporte,
no una advertencia al pie: una ocupación subestimada puede justificar la decisión
de no proteger un sitio.
